\documentclass[12pt]{article}
\usepackage[a4paper,margin=1in]{geometry}
\usepackage[utf8]{inputenc}
\usepackage[T1]{fontenc}
\usepackage{amsmath,amssymb}
\usepackage{textcomp}
\usepackage{parskip}
\begin{document}
\begin{center}
{\Large \textbf{2006 STEP I}}
\end{center}

\subsection*{Page 1}
\begin{flushleft}
STEP I, 2006 2\par
\bigskip\textbf{Section A: Pure Mathematics}\par
2 2\par
\bigskip\textbf{Question 1.} Find the integer, n, that satisfies n < 33 127 < (n + 1) . Find also a small integer m such\par
that (n + m) 33 127 is a perfect square. Hence express 33 127 in the form pq, where p and\par
-\par
q are integers greater than 1.\par
By considering the possible factorisations of 33 127, show that there are exactly two values of\par
m for which (n + m) 33 127 is a perfect square, and find the other value.\par
-\par
\bigskip\textbf{Question 2.} A small goat is tethered by a rope to a point at ground level on a side of a square barn which\par
stands in a large horizontal field of grass. The sides of the barn are of length 2a and the rope\par
is of length 4a. Let A be the area of the grass that the goat can graze. Prove that A 6 14\pia 2\par
and determine the minimum value of A.\par
\bigskip\textbf{Question 3.} In this question b, c, p and q are real numbers.\par
\textbf{(i)} By considering the graph y = x + bx + c show that c < 0 is a sufficient condition for\par
\[the equation x + bx + c = 0 to have distinct real roots. Determine whether c < 0 is a\]
necessary condition for the equation to have distinct real roots.\par
\textbf{(ii)} Determine necessary and sufficient conditions for the equation x + bx + c = 0 to have\par
distinct positive real roots.\par
\textbf{(iii)} What can be deduced about the number and the nature of the roots of the equation\par
\[x + px + q = 0 if p > 0 and q < 0?\]
\[What can be deduced if p < 0 and q < 0? You should consider the different cases that\]
3 2\par
arise according to the value of 4p + 27q .\par
\end{flushleft}
\subsection*{Page 2}
\begin{flushleft}
STEP I, 2006 3\par
\bigskip\textbf{Question 4.} By sketching on the same axes the graphs of y = sin x and y = x, show that, for x > 0:\par
\textbf{(i)} x > sin x ;\par
sin x\par
\textbf{(ii)} 1 for small x.\par
x \approx\par
A regular polygon has n sides, and perimeter P . Show that the area of the polygon is\par
P\par
.\par
\[\pi\]
4n tan\par
n\par
\gamma \delta\par
Show by differentiation (treating n as a continuous variable) that the area of the polygon\par
increases as n increases with P fixed.\par
Show also that, for large n, the ratio of the area of the polygon to the area of the smallest\par
circle which can be drawn around the polygon is approximately 1.\par
\bigskip\textbf{Question 5.} (i) Use the substitution u = 2x + 1 to show that, for x > 4,\par
\[3 \sqrt\{2x\} + 1 3\]
\[dx = ln - + K ,\]
\[Z (x 4)\sqrt\{2x\} + 1 \lambda \sqrt\{2x\} + 1 + 3 \mu\]
-\par
where K is a constant.\par
ln 8\par
2 7 2\par
\textbf{(ii)} Show that dx = + ln .\par
\[Zln 3 ex\sqrt\{ex\} + 1 12 3\]
2 2\par
\bigskip\textbf{Question 6.} (i) Show that, if (a , b) is any point on the curve x 2y = 1, then (3a + 4b , 2a + 3b) also\par
-\par
lies on the curve.\par
\textbf{(ii)} Determine the smallest positive integers M and N such that, if (a , b) is any point on\par
2 2\par
\[the curve M x N y = 1, then (5a + 6b , 4a + 5b) also lies on the curve.\]
-\par
2 2\par
\textbf{(iii)} Given that the point (a , b) lies on the curve x 3y = 1 , find positive integers P , Q,\par
-\par
R and S such that the point (P a + Qb , Ra + Sb) also lies on the curve.\par
\end{flushleft}
\subsection*{Page 3}
\begin{flushleft}
STEP I, 2006 4\par
\bigskip\textbf{Question 7.} (i) Sketch on the same axes the functions cosec x and 2x/\pi, for 0 < x < \pi . Deduce that\par
\[the equation x sin x = \pi/2 has exactly two roots in the interval 0 < x < \pi .\]
Show that\par
\[\pi \pi 3\pi 2\]
\[x sin x dx = 2 sin α + α\pi \pi 2α cos α 1\]
\[Z \pi/2 \{ - 2 \{ 4 - - - -\]
\{ \{\par
\{ \{\par
where α is the larger of the roots referred to above.\par
\textbf{(ii)} Show that the region bounded by the positive x-axis, the y-axis and the curve\par
\[y = ex 1 1\]
\{| - | - \{\par
\{ \{\par
\{ \{\par
has area ln 4 1.\par
-\par
8 Note that the volume of a tetrahedron is equal to the area of the base the height.\par
\[\times \times\]
The points O, A, B and C have coordinates (0, 0, 0), (a, 0, 0), (0, b, 0) and (0, 0, c), respectively,\par
where a, b and c are positive.\par
\textbf{(i)} Find, in terms of a, b and c, the volume of the tetrahedron OABC.\par
\textbf{(ii)} Let angle ACB = θ. Show that\par
\[cos θ =\]
\[(a^2 + c^2)(b^2 + c^2)\]
q\par
and find, in terms of a, b and c, the area of triangle ABC.\par
Hence show that d, the perpendicular distance of the origin from the triangle ABC, satisfies\par
1 1 1 1\par
\[= + + .\]
\[d^2 a^2 b^2 c^2\]
\end{flushleft}
\subsection*{Page 4}
\begin{flushleft}
STEP I, 2006 5\par
\bigskip\textbf{Section B: Mechanics}\par
\bigskip\textbf{Question 9.} A block of mass 4 kg is at rest on a smooth, horizontal table. A smooth pulley P is fixed to\par
one edge of the table and a smooth pulley Q is fixed to the opposite edge. The two pulleys\par
and the block lie in a straight line.\par
Two horizontal strings are attached to the block. One string runs over pulley P ; a particle of\par
mass x kg hangs at the end of this string. The other string runs over pulley Q; a particle of\par
\[mass y kg hangs at the end of this string, where x > y and x + y = 6 .\]
The system is released from rest with the strings taut. When the 4 kg block has moved a\par
distance d, the string connecting it to the particle of mass x kg is cut. Show that the time\par
taken by the block from the start of the motion until it first returns to rest (assuming that it\par
does not reach the edge of the table) is d/(5g) f (y), where\par
p\par
10 4\par
\[f (y) = + 1 + 6 2y.\]
\[\sqrt\{6\} 2y \lambda y \mu -\]
- p\par
\[Calculate the value of y for which f (y) = 0.\]
\bigskip\textbf{Question 10.} A particle P is projected in the x-y plane, where the y-axis is vertical and the x-axis is\par
\textasciicircum{}\circ\par
horizontal. The particle is projected with speed V from the origin at an angle of 45 above\par
the positive x-axis. Determine the equation of the trajectory of P .\par
The point of projection (the origin) is on the floor of a barn. The roof of the barn is given\par
\[by the equation y = x tan α + b , where b > 0 and α is an acute angle. Show that, if the\]
particle just touches the roof, then V ( 1 + tan α) = 2\sqrtbg; you should justify the choice of\par
- -\par
the negative root. If this condition is satisfied, find, in terms of α, V and g, the time after\par
projection at which touching takes place.\par
A particle Q can slide along a smooth rail fixed, in the x-y plane, to the under-side of the\par
roof. It is projected from the point (0, b) with speed U at the same time as P is projected\par
from the origin. Given that the particles just touch in the course of their motions, show that\par
\[2\sqrt\{2\} U cos α = V 2 + sin α cos α sin 2 α).\]
-\par
\end{flushleft}
\subsection*{Page 5}
\begin{flushleft}
STEP I, 2006 6\par
\[11 Particles A^1, A^2, A^3, . . ., A\]
n\par
\[(where n > 2) lie at rest in that order in a smooth straight\]
\[horizontal trough. The mass of A n-1 is m and the mass of A n is λm, where λ > 1. Another\]
particle, A0, of mass m, slides along the trough with speed u towards the particles and collides\par
\[with A^1. Momentum and energy are conserved in all collisions.\]
\textbf{(i)} Show that it is not possible for there to be exactly one particle moving after all collisions\par
have taken place.\par
\textbf{(ii)} Show that it is not possible for A n-1 and A n to be the only particles moving after all\par
collisions have taken place.\par
\textbf{(iii)} Show that it is not possible for A n-2, A n-1 and A n to be the only particles moving after\par
all collisions have taken place.\par
\textbf{(iv)} Given that there are exactly two particles moving after all collisions have taken place,\par
find the speeds of these particles in terms of u and λ.\par
\end{flushleft}
\subsection*{Page 6}
\begin{flushleft}
STEP I, 2006 7\par
\bigskip\textbf{Section C: Probability and Statistics}\par
12 Oxtown and Camville are connected by three roads, which are at risk of being blocked by\par
flooding. On two of the three roads there are two sections which may be blocked. On the\par
third road there is only one section which may be blocked. The probability that each section\par
is blocked is p. Each section is blocked independently of the other four sections. Show that\par
3 2\par
the probability that Oxtown is cut off from Camville is p (2 p) .\par
-\par
I want to travel from Oxtown to Camville. I choose one of the three roads at random and find\par
that my road is not blocked. Find the probability that I would not have reached Camville\par
if I had chosen either of the other two roads. You should factorise your answer as fully as\par
possible.\par
Comment briefly on the value of this probability in the limit p 1.\par
\to\par
\bigskip\textbf{Question 13.} A very generous shop-owner is hiding small diamonds in chocolate bars. Each diamond is\par
hidden independently of any other diamond, and on average there is one diamond per kilogram\par
of chocolate.\par
\textbf{(i)} I go to the shop and roll a fair six-sided die once. I decide that if I roll a score of N , I will\par
buy 100N grams of chocolate. Show that the probability that I will have no diamonds\par
is\par
e\par
-0.1\par
1 e\par
-0.6\par
-\par
6 \lambda 1 e-0.1 \mu\par
-\par
Show also that the expected number of diamonds I find is 0.35.\par
\textbf{(ii)} Instead, I decide to roll a fair six-sided die repeatedly until I score a 6. If I roll my first\par
6 on my T th throw, I will buy 100T grams of chocolate. Show that the probability that\par
I will have no diamonds is\par
e\par
-0.1\par
6 5e-0.1\par
-\par
Calculate also the expected number of diamonds that I find. (You may find it useful to\par
-2\par
consider the the binomial expansion of (1 x) .)\par
-\par
\bigskip\textbf{Question 14.} (i) A bag of sweets contains one red sweet and n blue sweets. I take a sweet from the bag,\par
note its colour, return it to the bag, then shake the bag. I repeat this until the sweet I\par
take is the red one. Find an expression for the probability that I take the red sweet on\par
the rth attempt. What value of n maximises this probability?\par
\textbf{(ii)} Instead, I take sweets from the bag, without replacing them in the bag, until I take the\par
red sweet. Find an expression for the probability that I take the red sweet on the rth\par
attempt. What value of n maximises this probability?\par
\end{flushleft}
\end{document}
